% ============================================================
% f2015ahl_q07_views.tex —— 2015 安徽卷（理）第 7 题 三视图与尺寸标注
%
% 【来源】工作档 .work/phase3-b13/source/2015安徽理.pdf（sha256 402eca00…ce7b0）
%          第 1 页（页脚 931）栏 1；坐标数据取自实测清单
%          .work/phase3-b13/figure-inventory.md §F01（要素表 P031–P075、箭头表、标签表）。
% 【坐标口径】真页 pt（PDF 页面坐标，原点页左下角，1190.55 × 841.89 bp）。
%          画布 1 单位 = 1 TeX pt，图内坐标直接写字面值（清单已给真页 pt）。
%          清单 §0.3 本卷实测换算因子 a_x = 0.999622795、a_y = 0.999621571，
%          与 1 的偏差 3.8e-4 系渲染器取整所致（非页面缩放），故不再二次缩放。
%          注：1 TeX pt = 0.996264 bp，故本图绝对尺寸比源件小 0.375%（项目定型口径）。
% 【线宽】0.3980 pt：3 个三角形三边、2 条高、12 条尺寸线本体、10 条竖刻度、
%          14 条 √2 引线（本体+刻度）；0.4200 pt：4 处「√2」根号上横线。
%          注：清单 §F01.2(c) 小标题写「尺寸线本体 lw = 0.42 pt」，与其 §F01.2 末尾
%          「线宽分档」表（0.398 × 45 + 0.42 × 4）及源件实测（P007…P029 全为 0.3980）
%          矛盾 —— 此处按源件实测与分档表取 0.3980，仅 4 条根号上横线为 0.4200。
% 【照源复刻的存疑项】
%   1) 源件中俯视图右侧两条 √2 引线刻度被重绘一次（P070/P071 与 P056/P057 端点逐位
%      相同，清单 FI-S1）—— 几何重合，此处只画一次。
%   2) 4 处「√2」的根号上横线（P072–P075，lw 0.42）不另画路径：本图用 `$\sqrt{2}$`
%      排版，横线由 TeX 根号构造内生（与源件同构，源件亦为 CMSY10 根号 + 独立横线）。
%   3) 三视图形「底宽 2、高 1」与侧视图「1 × 1」的几何基准不统一（FI-S15），
%      照源复刻不改。
%   4) 标签按【源墨迹实测中心】放置；清单给的是 pdftohtml em 框中心（西文含 descend，
%      中心偏低约 1.3–1.5 pt；CMSY10 根号框另有上偏），故坐标 = 清单值 + 墨迹修正。
%   5) 「正(主)视图」等中文标签用半角括号直排（源件括号为 LMRoman10 半角），
%      中西文间距由 ctex 默认 CJKecglue 生成，与源件的间距可能略有差异。
% ============================================================
% 【自验】（探针 .work/phase3-b13/probe/q07-probe.tex，latexmk -pdfxe）
%   编译：0 error、0 Overfull、0 Underfull、0 Missing character。
%   字面量反算双层核对（A 手工誊录清单字面量 + B 源 PDF 逐命令复解析，含两条
%   曲线的全部 598 段控制点）：最大残差 0.0005 pt（阈值 0.05）。
%   600 dpi 墨迹核对（源件页 vs 探针对位后）：几何 345 点采样 max 0.145 pt / rms 0.109 pt，标签墨迹质心 max 0.127 pt；
%   并排与红蓝重合裁片见 probe/crops/q07-side.png、q07-overlay.png。
\begin{tikzpicture}[x=1pt, y=1pt,
    every node/.style={font=\fontsize{10.5pt}{12.6pt}\selectfont, inner sep=0pt}]
  \begin{scope}[line width=0.398pt]
    % ---------- 三角形与高（清单 §F01.2(a)）----------
    \draw (135.701,215.511) -- (223.009,215.511) -- (179.355,259.165) -- cycle;  % P044 正视图
    \draw (179.355,215.511) -- (179.355,259.165);                                % P043 正视图高
    \draw (135.701,150.030) -- (223.009,150.030) -- (179.355,106.376) -- cycle;  % P041 俯视图
    \draw (179.355,150.030) -- (179.355,106.376);                                % P042 俯视图高
    \draw (266.663,215.511) -- (310.317,215.511) -- (266.663,259.165) -- cycle;  % P045 侧视图
    % ---------- 尺寸刻度线（§F01.2(b)，10 条）----------
    \draw (135.701,154.396) -- (135.701,163.127);   % P036 俯视图行左
    \draw (179.355,154.396) -- (179.355,163.127);   % P037 俯视图行中
    \draw (223.009,154.396) -- (223.009,163.127);   % P038 俯视图行右
    \draw (135.701,202.414) -- (135.701,211.145);   % P033 正视图行左
    \draw (179.355,202.414) -- (179.355,211.145);   % P034 正视图行中
    \draw (223.009,202.414) -- (223.009,211.145);   % P035 正视图行右
    \draw (266.663,202.414) -- (266.663,211.145);   % P039 侧视图行左
    \draw (310.317,202.414) -- (310.317,211.145);   % P040 侧视图行右
    \draw (253.566,215.511) -- (262.297,215.511);   % P031 侧视图竖标下
    \draw (253.566,259.165) -- (262.297,259.165);   % P032 侧视图竖标上
    % ---------- 尺寸线本体（§F01.2(c)，12 条）----------
    \draw (139.2872,158.7605) -- (148.7972,158.7605);   % P007 俯视图左段
    \draw (166.2585,158.7605) -- (175.7685,158.7605);   % P009 俯视图左段尾
    \draw (182.9411,158.7605) -- (192.4511,158.7605);   % P011 俯视图右段首
    \draw (209.9124,158.7605) -- (219.4224,158.7605);   % P013 俯视图右段
    \draw (139.2872,206.7800) -- (148.7972,206.7800);   % P015 正视图左段
    \draw (166.2585,206.7800) -- (175.7685,206.7800);   % P017 正视图左段尾
    \draw (182.9411,206.7800) -- (192.4511,206.7800);   % P019 正视图右段首
    \draw (209.9124,206.7800) -- (219.4224,206.7800);   % P021 正视图右段
    \draw (270.2489,206.7800) -- (279.7589,206.7800);   % P023 侧视图横标左
    \draw (297.2202,206.7800) -- (306.7302,206.7800);   % P025 侧视图横标右
    \draw (257.9320,219.0970) -- (257.9320,228.6070);   % P027 侧视图竖标上
    \draw (257.9320,246.0680) -- (257.9320,255.5780);   % P029 侧视图竖标下
    % ---------- √2 引线与刻度（§F01.2(c)，P070/P071 与 P056/P057 重合故略）----------
    \draw (133.518,147.847) -- (126.533,140.862);   % P046 俯视图左刻度
    \draw (177.172,104.193) -- (170.187,97.208);    % P047 俯视图左刻度（下）
    \draw (143.122,131.258) -- (132.562,141.818);   % P048 俯视图左引线本体
    \draw (160.583,113.797) -- (171.143,103.237);   % P050 俯视图左引线本体
    \draw (215.588,131.258) -- (226.148,141.818);   % P052 俯视图右引线本体
    \draw (198.126,113.797) -- (187.566,103.237);   % P054 俯视图右引线本体
    \draw (225.191,147.847) -- (232.177,140.862);   % P056 俯视图右刻度
    \draw (181.538,104.193) -- (188.523,97.208);    % P057 俯视图右刻度（下）
    \draw (133.518,217.694) -- (126.533,224.679);   % P058 正视图左刻度
    \draw (177.172,261.347) -- (170.187,268.332);   % P059 正视图左刻度（上）
    \draw (225.191,217.694) -- (232.177,224.679);   % P060 正视图右刻度
    \draw (181.538,261.347) -- (188.523,268.332);   % P061 正视图右刻度（上）
    \draw (143.122,234.282) -- (132.562,223.722);   % P062 正视图左引线本体
    \draw (160.583,251.744) -- (171.143,262.303);   % P064 正视图左引线本体
    \draw (215.588,234.282) -- (226.148,223.722);   % P066 正视图右引线本体
    \draw (198.126,251.744) -- (187.566,262.303);   % P068 正视图右引线本体
    % ---------- 20 个填充箭头（尖端坐标取清单；Triangle 长 3.985 / 宽 2.988）----------
    \draw[{Triangle[length=3.985pt,width=2.988pt]}-] (135.7009,158.7605) -- (139.6859,158.7605);  % P008 ← 俯视图左端
    \draw[{Triangle[length=3.985pt,width=2.988pt]}-] (179.3547,158.7605) -- (175.3697,158.7605);  % P010 → 俯视图中点
    \draw[{Triangle[length=3.985pt,width=2.988pt]}-] (179.3548,158.7605) -- (183.3398,158.7605);  % P012 ← 俯视图中点
    \draw[{Triangle[length=3.985pt,width=2.988pt]}-] (223.0086,158.7605) -- (219.0236,158.7605);  % P014 → 俯视图右端
    \draw[{Triangle[length=3.985pt,width=2.988pt]}-] (135.7009,206.7800) -- (139.6859,206.7800);  % P016 ← 正视图左端
    \draw[{Triangle[length=3.985pt,width=2.988pt]}-] (179.3547,206.7800) -- (175.3697,206.7800);  % P018 → 正视图中点
    \draw[{Triangle[length=3.985pt,width=2.988pt]}-] (179.3548,206.7800) -- (183.3398,206.7800);  % P020 ← 正视图中点
    \draw[{Triangle[length=3.985pt,width=2.988pt]}-] (223.0086,206.7800) -- (219.0236,206.7800);  % P022 → 正视图右端
    \draw[{Triangle[length=3.985pt,width=2.988pt]}-] (266.6626,206.7800) -- (270.6476,206.7800);  % P024 ← 侧视图左端
    \draw[{Triangle[length=3.985pt,width=2.988pt]}-] (310.3164,206.7800) -- (306.3314,206.7800);  % P026 → 侧视图右端
    \draw[{Triangle[length=3.985pt,width=2.988pt]}-] (257.9319,215.5107) -- (257.9319,219.4957);  % P028 ↓ 侧视图竖标下端
    \draw[{Triangle[length=3.985pt,width=2.988pt]}-] (257.9319,259.1645) -- (257.9319,255.1795);  % P030 ↑ 侧视图竖标上端
    \draw[{Triangle[length=3.985pt,width=2.988pt]}-] (130.0257,144.3546) -- (132.8436,141.5368);  % P049 ↖ 俯视图左 √2 上端
    \draw[{Triangle[length=3.985pt,width=2.988pt]}-] (173.6795,100.7008) -- (170.8616,103.5186);  % P051 ↘ 俯视图左 √2 下端
    \draw[{Triangle[length=3.985pt,width=2.988pt]}-] (228.6838,144.3546) -- (225.8660,141.5368);  % P053 ↗ 俯视图右 √2 上端
    \draw[{Triangle[length=3.985pt,width=2.988pt]}-] (185.0300,100.7008) -- (187.8479,103.5186);  % P055 ↙ 俯视图右 √2 下端
    \draw[{Triangle[length=3.985pt,width=2.988pt]}-] (130.0257,221.1859) -- (132.8436,224.0037);  % P063 ↖ 正视图左 √2 上端
    \draw[{Triangle[length=3.985pt,width=2.988pt]}-] (173.6795,264.8397) -- (170.8616,262.0219);  % P065 ↗ 正视图左 √2 下端
    \draw[{Triangle[length=3.985pt,width=2.988pt]}-] (228.6838,221.1859) -- (225.8660,224.0037);  % P067 ↗ 正视图右 √2 上端
    \draw[{Triangle[length=3.985pt,width=2.988pt]}-] (185.0300,264.8397) -- (187.8479,262.0219);  % P069 ↙ 正视图右 √2 下端
  \end{scope}
  % ---------- 文字标签（源墨迹实测中心；行末为清单 em 框中心，供追溯）----------
  \node at (146.621,121.640) {$\sqrt{2}$};   % L q07_sqrt2_tl  √(143.733,129.419) 2(151.403,119.412)
  \node at (206.007,121.076) {$\sqrt{2}$};   % L q07_sqrt2_tr  √(203.094,129.419) 2(210.097,119.412)
  \node at (157.485,206.699) {$1$};          % L q07_1_fl      (157.406,205.469)
  \node at (201.120,206.699) {$1$};          % L q07_1_fr      (201.426,205.469)
  \node at (157.492,158.792) {$1$};          % L q07_1_tl      (157.406,157.437)
  \node at (201.137,158.742) {$1$};          % L q07_1_tr      (201.426,157.437)
  \node at (288.506,206.699) {$1$};          % L q07_1_side_h  (288.800,205.469)
  \node at (257.938,237.218) {$1$};          % L q07_1_side_v  (258.119,236.156)
  \node at (179.108,184.927) {正(主)视图};    % L q07_front     核字框并集 x[150.736,208.096] y[179.452,190.793]
  \node at (288.329,184.825) {侧(左)视图};    % L q07_side      核字框并集 x[260.120,317.480] y[179.452,190.793]
  \node at (179.389,84.540) {俯视图};         % L q07_top       (179.083,84.723)
  \node at (146.725,243.978) {$\sqrt{2}$};   % L q07_sqrt2_fl  √(143.733,251.500) 2(151.403,242.160)
  \node at (206.132,244.216) {$\sqrt{2}$};   % L q07_sqrt2_fr  √(203.094,251.500) 2(210.097,242.160)
\end{tikzpicture}
